% ============================================================
% exp_security.tex - 安全方向工作经历
% 适用岗位：安全开发、API安全、WAF、HIDS 等
% ============================================================

\section{工作经验}
\vspace{0.2cm}

\datedsubsection{闪捷信息科技有限公司}{安全技术专家}{2022.4-至今}

主要负责API安全网关、API审计类产品的设计和核心开发工作,以及项目管理工作。主要功能包括：Http代理，账号发现，敏感数据识别，脱敏，水印，访问控制，接口梳理，风险与脆弱性发现等。

\vspace{0.2cm}

\textit{工程开发(c/c++)}
\begin{itemize}
    \setstretch{1.2}
    \item 负责整个网关项目的技术架构，极致性能优化，疑难阻塞问题攻坚等;
    \item 基于Nginx二次开发，Module开发，较之原有代码量减少了一半，容错性大幅提升;
    \item 网关核心模块开发访问控制，账号发现，Agent, 分布式部署等功能设计与开发;
\end{itemize}

\textit{项目管理}
\begin{itemize}
\item 网关项目PM，主线开发，支撑POC测试，项目交付。敏捷开发, 持续集成，快速迭代;
\end{itemize}

\textit{产品及市场支持}
\begin{itemize}
    \setstretch{1.2}
    \item 前入一线，理解客户需求，从风险视角和安全体系建设的视角宣导产品价值，参与客户数据安全方案制定，收集在稳定性、稳定性担忧，功能等方面问题形成产品需求;
    \item 输出产品需求，扩展功能，打造产品亮点, 其中包括：风险规则、告警抑制、稳定性方案、交互体验等;
\end{itemize}

\vspace{0.2cm}
\datedsubsection{玩物得志(杭州装点文化创意有限公司)}{安全专家, 团队leader}{2019.8-2022.3}
安全风控团队负责人，从0-1组织搭建玩物安全风控体系。主要负责安全风控核心开发，安全风控运营，等保合规理等工作。

\textit{风控体系建设}
\begin{itemize}
    \setstretch{1.2}
  \item 基础工程: 风控网关，业务风控，内容风控，数据分析，管理台;
  \item 风控业务: 特征提取=>数据分析=>制定整体方案=>推动落地；黑灰产对抗，恶意注册，人机对抗，反爬虫，薅羊毛，反洗钱，内容反垃圾，私域导流，IM诈骗、监管合规等;
\end{itemize}

\textit{安全体系建设}
\begin{itemize}
    \setstretch{1.3}
  \item 主机安全： 安全意识培训，危险入侵报警维护;
  \item 安全审计： 网络安全审计，VPN，收敛风险暴露面；数据安全审计，登录认证和权限认证;
  \item 安全开发： 等保2.0，设备指纹，接口保护，客户端安全加固，反爬虫，渗透测试，人机对抗等;
\end{itemize}

\textit{文本反垃圾系统}

基于使用三方产品费用过高的问题，自建文本反垃圾系统，快速灵活响应需求，接入场景包括IM，评论，社区，昵称，商品描述等。
\begin{itemize}
    \setstretch{1.2}
    \item 核心基于DFA算法的敏感词匹配，同时支持正则，自建规则的内容匹配;
    \item 管理台，前后端各种接口;
\end{itemize}

\vspace{0.2cm}
\datedsubsection{蘑菇街(美丽联合集团)}{资深开发工程师}{2015.5-2019.8}
作为核心人员，主要从事IM系统，网络安全，HIDS，WAF等开发维护工作；见证并参与了自研WAF从0-1，以及整个mogu安全逐步建成一个完备的安全体系的过程。

\vspace{0.2cm}
\textit{安全开发 (C++/Golang)}
\begin{itemize}
    \setstretch{1.2}
    \item 作为核心开发人员自研Http层网关和Waf开发；接入全站的流量，单机性能在3w+/s; 10ms超时率在 10万分之5以下；
    \item 主机入侵检测系统开发和维护工作；支持命令监控，文件监控，漏洞扫描，堡垒机安装，安全基线扫描等功能；
    \item 基于ipset和iptables的主机ACL系统开发和爬虫分析等其他开发类工作；
\end{itemize}

\textit{IM后端 (C++)}
\begin{itemize}
    \setstretch{1.2}
\item IM2.0 重构版本开发，主要负责接入服务器、登录、会话、长链推送等业务模块；
\item 独立开发完成 IM 客户端 SDK 网络库，后推广为全公司的通用长连接网络库；
\end{itemize}

\textit{信息安全}
\begin{itemize}
    \setstretch{1.2}
    \item 主要负责主机安全扫描，中间件安全扫描同时推动业务方进行安全整改；
    \item 对安全代码进行白盒扫描，推动整改相关安全问题；漏洞验证与公司内部组织的渗透测试；
\end{itemize}

\vspace{0.2cm}
\datedsubsection{北京千方科技集团有限公司}{研发经理}{2012.4-2015.5}
全球最大车联网接入平台；车辆GPS监管、监控、交通部货代平台的大数据接入。

\vspace{0.2cm}
\textit{项目管理}
\begin{itemize}
    \setstretch{1.2}
    \item 负责整个服务端通讯团队的项目进度把控，发起立项申请，项目时间规划，任务安排及资源协调等工作。开发过程中采用敏捷开发方式，多版本并行开发，注重开发节奏控制；
\end{itemize}

\textit{架构 \& 编码 (c++)}
\begin{itemize}
    \setstretch{1.2}
    \item 通过多个版本迭代改造，在主要核心服务之间实现了服务自发现，连接自调度，自动负载均衡，服务热扩展；
    \item 平衡性能和数据一致性的矛盾关系，通过静态数据广播，动态数据扩散，缓存管理等方式，使得由于集群节点之间数据一致性问题导致的车机上线异常问题由1例/1周降为小于1例/3月；
    \item 基于事件驱动模型的车机数据透传服务获得国家发明专利；
\end{itemize}

\vspace{0.2cm}
\datedsubsection{华为技术有限公司 | 北研所}{开发工程师}{2011.5-2012.4}
    \setstretch{1.2}
华为电信级路由器VRPV8软件支撑平台的开发工作。平台基于Linux OS裁剪，在面对时间紧、任务重、工作量大的情况下，在保自身开发进度的同时，多次协助定位疑难BUG，有力保证了项目顺利过点。

\vspace{0.2cm}
\textit{开发技术：C、Linux，网络编程，文件操作，敏捷开发，持续集成}

\vspace{0.2cm}
\datedsubsection{北京千方科技集团有限公司 | 集团研发}{开发工程师}{2009.11（实习）-2011.5}
    \setstretch{1.2}
交通部信息中心的全国车辆监管平台的研发工作。对接广州亚运会"两客一危"车辆GPS监控系统，独立承担后台多个系统的从设计，开发，测试上线的工作，覆盖从GPS终端网关接入，协议解析，分发，入库等整个流程，并多次赴实地运维开发工作，荣获优秀实习生，优秀员工。

\vspace{0.2cm}
\textit{开发技术：C++，TCP/UDP，Linux，网关，LVS，GPS通信协议，Mysql}
